% ============================================================================
% LEHRERHINWEISE: Verbos en -ar (Repaso) im Schulkontext
% Klasse 7 | Spanisch | Unidad 2 | DS 04
% ============================================================================

\documentclass[11pt,a4paper]{article}

\usepackage[utf8]{inputenc}
\usepackage[T1]{fontenc}
\usepackage[ngerman]{babel}
\usepackage{geometry}
\usepackage{fancyhdr}
\usepackage{xcolor}
\usepackage{enumitem}
\usepackage{tabularx}
\usepackage{array}
\usepackage{tcolorbox}
\usepackage{booktabs}
\usepackage{amssymb}

% ============================================================================
% LAYOUT
% ============================================================================
\geometry{left=2cm, right=2cm, top=2cm, bottom=2cm}
\definecolor{fach}{HTML}{c41e3a}
\definecolor{grau}{HTML}{666666}
\definecolor{hellgrau}{HTML}{f5f5f5}
\definecolor{basis}{HTML}{2a7a4b}
\definecolor{standard}{HTML}{2c5aa0}
\definecolor{erweiterung}{HTML}{7b1fa2}
\definecolor{spiel}{HTML}{b35c00}

\pagestyle{fancy}
\fancyhf{}
\fancyhead[L]{\small\textcolor{grau}{Spanisch Klasse 7 -- Unidad 2}}
\fancyhead[R]{\small\textcolor{grau}{LH-04 (Lehrerhinweise)}}
\fancyfoot[C]{\small\thepage}
\renewcommand{\headrulewidth}{0.4pt}

% Spanisch-Befehl
\newcommand{\es}[1]{\textbf{#1}}

% ============================================================================
% DOKUMENT
% ============================================================================
\begin{document}

% ----------------------------------------------------------------------------
% KOPFBEREICH
% ----------------------------------------------------------------------------
\begin{center}
{\Large\textcolor{fach}{\textbf{Lehrerhinweise}}}\\[0.3cm]
{\large Verbos en -ar (Repaso) im Schulkontext}\\[0.2cm]
{\normalsize Doppelstunde (90 Min.) | Qué pasa 1, Repaso Unidad 1+2}
\end{center}

\vspace{0.5cm}
\hrule
\vspace{0.5cm}

% ----------------------------------------------------------------------------
% KURZÜBERSICHT
% ----------------------------------------------------------------------------
\section*{Kurzübersicht}

\begin{tabular}{ll}
\textbf{Thema:} & -ar Verben wiederholen + Schulalltag beschreiben \\
\textbf{Lehrwerk:} & Qué pasa 1, Repaso Unidad 1+2 \\
\textbf{Materialien:} & AB-04, LB-04, ML-04, Verbkarten, Ball \\
\textbf{Fokus:} & Konjugation -ar festigen, Leseverstehen, Textproduktion \\
\textbf{Voraussetzung:} & hay, estar, de+art. (DS 01-03), ser, tener (Unidad 1) \\
\end{tabular}

\vspace{0.5cm}

% ----------------------------------------------------------------------------
% FEINZIELE
% ----------------------------------------------------------------------------
\section*{Feinziele}

\begin{enumerate}[label=\textbf{FZ\arabic*:}, leftmargin=2.5em]
    \item SuS \textbf{konjugieren} -ar Verben korrekt in allen Personen \textit{(AFB I)}
    \item SuS \textbf{verwenden} -ar Verben im Schulkontext \textit{(AFB II)}
    \item SuS \textbf{kombinieren} ser, tener, estar, hay mit -ar Verben in Texten \textit{(AFB II)}
    \item SuS \textbf{verstehen} einen einfachen Text über den Schulalltag \textit{(AFB I)}
\end{enumerate}

\vspace{0.5cm}

% ----------------------------------------------------------------------------
% STUNDENVERLAUF
% ----------------------------------------------------------------------------
\section*{Stundenverlauf (90 Min.)}

\begin{tabular}{|p{1cm}|p{2.2cm}|p{5.8cm}|p{1cm}|p{3cm}|}
\hline
\textbf{Min} & \textbf{Phase} & \textbf{Aktivität} & \textbf{SF} & \textbf{Material/Signal} \\
\hline
0--8 & Einstieg & Verb-Pantomime: L macht Verb vor, SuS rufen Infinitiv & PL & -- \\
\hline
8--15 & Repaso & LB: -ar Konjugation, Endungen markieren, 3 Beispiele & EA & LB-04, \newline\textbf{Signal:} Gong \\
\hline
15--28 & Übung 1 & AB Aufg. 1--2: Konjugationstabelle + Lückensätze & EA & AB-04, Timer 13' \\
\hline
28--42 & Leseverstehen & Text \es{Un día en el instituto} lesen, Fragen PA & EA$\rightarrow$PA & LB-04, \newline\textbf{Signal:} Partner \\
\hline
42--47 & Pause & Aufstehen, strecken & -- & Musik \\
\hline
47--60 & Übung 2 & AB Aufg. 3--4: Textverständnis + Sätze bilden & EA & AB-04, Timer 13' \\
\hline
60--70 & \textbf{Spiel} & \es{Cadena de verbos} mit Verbkarten + Ball & PL & Karten, Ball, \newline\textbf{Signal:} Musik aus \\
\hline
70--85 & Produktion & AB Aufg. 5: Mini-Text (3--4 Sätze) mit Satzstartern & EA & AB-04, Timer 15' \\
\hline
85--90 & Sicherung & 1--2 Freiwillige lesen Text vor, Applaus & PL & Applaus \\
\hline
\end{tabular}

\vspace{0.3cm}
\textbf{Zeitpuffer:} In Produktion integriert (15 Min)

\vspace{0.5cm}

% ----------------------------------------------------------------------------
% DIFFERENZIERUNG
% ----------------------------------------------------------------------------
\section*{Differenzierung}

\begin{tcolorbox}[colback=white, colframe=basis, title=\textcolor{white}{\textbf{[B] Basis}}]
\begin{itemize}[nosep]
    \item Aufg. 1: Konjugationstabelle mit Endungen vorgegeben (nur einsetzen)
    \item Aufg. 2: Lückensätze mit Verb-Infinitiv in Klammern
    \item Text: Mit Glossar (10 schwierige Wörter übersetzt) -- LB-04 ausschneiden!
    \item Aufg. 5: 3 Sätze, mit Satzstartern
    \item Spiel: Bekommt Verbkarte, muss nur konjugieren (nicht frei wählen)
\end{itemize}
\end{tcolorbox}

\begin{tcolorbox}[colback=white, colframe=standard, title=\textcolor{white}{\textbf{[S] Standard}}]
\begin{itemize}[nosep]
    \item Aufg. 1--5: Vollständig ohne Hilfe
    \item Text: Ohne Glossar
    \item Spiel: Normale Teilnahme
    \item Aufg. 5: 3--4 Sätze
\end{itemize}
\end{tcolorbox}

\begin{tcolorbox}[colback=white, colframe=erweiterung, title=\textcolor{white}{\textbf{[E] Erweiterung}}]
\begin{itemize}[nosep]
    \item Aufg. 1--4: Schnell fertig $\rightarrow$ Zusatz: Text weiterschreiben
    \item Spiel: Muss 2 Verben in einem Satz kombinieren
    \item Aufg. 5: 5--6 Sätze, alle Verbtypen (ser, tener, estar, hay, -ar)
\end{itemize}
\end{tcolorbox}

\vspace{0.5cm}

% ----------------------------------------------------------------------------
% SPIELANLEITUNG
% ----------------------------------------------------------------------------
\section*{Spiel: Cadena de verbos (Verb-Kette)}

\begin{tcolorbox}[colback=white, colframe=spiel, title=\textcolor{white}{\textbf{Spielablauf (10 Min.)}}]

\textbf{Material:} 12 Verbkarten (laminiert, A6), weicher Ball

\vspace{0.3cm}

\textbf{Verbkarten vorbereiten:}
hablar, estudiar, escuchar, mirar, buscar, tomar, descansar, charlar, esperar, ser, tener, estar

\vspace{0.3cm}

\textbf{Vorbereitung (2 Min):}
\begin{enumerate}[nosep]
    \item Alle stehen im Kreis
    \item Verbkarten werden an schwächere SuS verteilt (Hilfe!)
    \item Ball bei L
\end{enumerate}

\vspace{0.3cm}

\textbf{Ablauf:}
\begin{enumerate}[nosep]
    \item L beginnt: \es{Yo hablo español.} $\rightarrow$ wirft Ball
    \item Fänger muss:
    \begin{itemize}[nosep]
        \item \textbf{Option A:} Anderes Verb + Objekt (\es{Tú estudias matemáticas.})
        \item \textbf{Option B:} Gleiches Verb, andere Person (\es{María habla inglés.})
        \item \textbf{Option C (mit Karte):} Verb von Karte konjugieren
    \end{itemize}
    \item Wer länger als 5 Sekunden braucht $\rightarrow$ darf Karte ziehen
    \item Wer Fehler macht $\rightarrow$ einmal aussetzen (nicht raus!)
    \item Nach 8 Min: Alle noch Stehenden = Gewinner
\end{enumerate}

\vspace{0.3cm}

\textbf{Regeln (an Tafel):}
\begin{enumerate}[nosep]
    \item Ball fangen $\rightarrow$ Satz sagen
    \item Max. 5 Sekunden!
    \item Hilfe? $\rightarrow$ Karte ziehen
    \item Fehler? $\rightarrow$ 1× aussetzen
    \item Immer freundlich!
\end{enumerate}

\vspace{0.3cm}

\textbf{Notfallplan (falls zu schwer):}

$\rightarrow$ Umwandeln in \glqq Konjugations-Staffel\grqq:
2 Teams, L sagt Verb + Person, Erster der richtig konjugiert bekommt Punkt.

\end{tcolorbox}

\vspace{0.5cm}

% ----------------------------------------------------------------------------
% ÜBERGANGSSIGNALE
% ----------------------------------------------------------------------------
\section*{Übergangssignale}

\begin{tabular}{|l|l|l|}
\hline
\textbf{Von $\rightarrow$ Nach} & \textbf{Signal} & \textbf{Ansage} \\
\hline
Einstieg $\rightarrow$ Repaso & Gong & \es{Abrimos el LB} \\
\hline
Repaso $\rightarrow$ Übung 1 & Timer Start & \es{AB, 13 minutos} \\
\hline
Übung 1 $\rightarrow$ Leseverstehen & \glqq Partner!\grqq & \es{Leemos el texto con pareja} \\
\hline
Text $\rightarrow$ Pause & Musik & \es{Pausa, 5 minutos} \\
\hline
Pause $\rightarrow$ Übung 2 & Musik aus & \es{Sentaos, AB Aufgabe 3} \\
\hline
Übung 2 $\rightarrow$ Spiel & \es{¡Levantaos!} & \es{Círculo, juego de verbos} \\
\hline
Spiel $\rightarrow$ Produktion & Ball weg & \es{Sentaos, escribimos 3-4 frases} \\
\hline
Produktion $\rightarrow$ Sicherung & Timer Ende & \es{¿Quién lee su texto?} \\
\hline
\end{tabular}

\vspace{0.5cm}

% ----------------------------------------------------------------------------
% MATERIALCHECKLISTE
% ----------------------------------------------------------------------------
\section*{Materialcheckliste}

\textbf{Vom Lehrer vorzubereiten:}
\begin{itemize}[nosep]
    \item[$\square$] \textbf{Verbkarten} (12 Stück, A6, laminiert):
    \begin{itemize}[nosep]
        \item hablar, estudiar, escuchar, mirar, buscar, tomar
        \item descansar, charlar, esperar, ser, tener, estar
    \end{itemize}
    \item[$\square$] \textbf{Weicher Ball} (Softball/Tennisball) für Verb-Kette
    \item[$\square$] \textbf{Timer} (Handy, sichtbar für alle)
    \item[$\square$] \textbf{Glossar-Kopien} für [B]-SuS (10 Stück) -- aus LB-04 ausschneiden
    \item[$\square$] \textbf{Satzstarter-Kärtchen} für Produktion (Klasse hat diese bereits)
\end{itemize}

\vspace{0.3cm}

\textbf{Kopien:}
\begin{itemize}[nosep]
    \item[$\square$] AB-04 (Klassensatz)
    \item[$\square$] LB-04 (Klassensatz)
    \item[$\square$] Lehrwerk Qué pasa 1 (Repaso)
\end{itemize}

\vspace{0.3cm}

\textbf{Von SuS mitzubringen:}
\begin{itemize}[nosep]
    \item LB (Lernblatt von heute)
    \item AB aus DS 03
\end{itemize}

\vspace{0.5cm}

% ----------------------------------------------------------------------------
% HÄUFIGE FEHLER
% ----------------------------------------------------------------------------
\section*{Häufige Fehler und Reaktionen}

\begin{tabular}{|p{3.5cm}|p{3.5cm}|p{6cm}|}
\hline
\textbf{Fehler} & \textbf{Ursache} & \textbf{Reaktion} \\
\hline
*\es{yo estudio} falsch betont & Akzentregel unklar & Stamm betonen: esTUdio, nicht estuDIO \\
\hline
*\es{nosotros hablan} & Person verwechselt & nosotros = -amos! (hablan = ellos) \\
\hline
*\es{él estudiar} & Nicht konjugiert & -ar weg, Endung dran: él estudia \\
\hline
*\es{vosotros estudiais} & Akzent vergessen & estudiÁIS -- Akzent auf dem i! \\
\hline
*\es{escuchar el profesor} & \es{a} fehlt & \es{escuchar AL profesor} (a + el = al) \\
\hline
\end{tabular}

\vspace{0.5cm}

% ----------------------------------------------------------------------------
% SATZSTARTER
% ----------------------------------------------------------------------------
\section*{Satzstarter für Produktion (AB Aufg. 5)}

\begin{tcolorbox}[colback=hellgrau, colframe=grau]
\textbf{An Tafel oder Folie sichtbar:}

\vspace{0.3cm}
\begin{itemize}[nosep]
    \item \es{Soy alumno/alumna en...}
    \item \es{Tengo ... años.}
    \item \es{En mi clase hay...}
    \item \es{Mi profesor/a de español se llama...}
    \item \es{Estudio... / Hablo... / Escucho...}
    \item \es{En el recreo descanso y...}
    \item \es{Me gusta...}
\end{itemize}
\end{tcolorbox}

\vspace{0.5cm}

% ----------------------------------------------------------------------------
% AB-LÖSUNGEN
% ----------------------------------------------------------------------------
\section*{AB-Lösungen (Kurzform)}

\textbf{Aufg. 1 -- Konjugationstabelle:}

\begin{tabular}{l|lll}
 & estudiar & escuchar & mirar \\
\hline
yo & estudio & escucho & miro \\
tú & estudias & escuchas & miras \\
él/ella & estudia & escucha & mira \\
nosotros & estudiamos & escuchamos & miramos \\
ellos & estudian & escuchan & miran \\
\end{tabular}

\vspace{0.3cm}

\textbf{Aufg. 2 -- Lücken:}
a) hablamos, b) miran, c) estudio, d) busca, e) escucháis, f) descansan

\vspace{0.3cm}

\textbf{Aufg. 3 -- Leseverständnis:}
\begin{itemize}[nosep]
    \item a) Hay veinticinco alumnos.
    \item b) Habla español y mira la pizarra.
    \item c) Descansan, charlan y buscan un banco.
\end{itemize}

\vspace{0.3cm}

\textbf{Aufg. 4 -- Sätze bilden:} Individuelle Lösungen, z.B.:
\begin{itemize}[nosep]
    \item Los alumnos escuchan al profesor en la clase.
    \item Nosotros estudiamos español y miramos la pizarra.
    \item María busca el libro. El libro está en la mochila.
\end{itemize}

\vspace{0.3cm}

\textbf{Aufg. 5 -- Mini-Text:} Individuelle Lösungen. Kriterien:
\begin{itemize}[nosep]
    \item Mind. 5 Sätze
    \item Mind. 3 verschiedene -ar Verben
    \item Mind. 1× ser/tener/estar/hay
    \item Korrekter Schulkontext
\end{itemize}

\vspace{0.5cm}

\end{document}
